\section{The route through one's own case}
\label{sec:19.3}
The construction most readers will propose first is not caregiving but a body: give the machine needs of its own, and let it recognize need in others by its own case. It is the older idea, it is Adam Smith's first step, and it is what the boundary model becomes if it is asked to carry the conscience. Each link of that route fails in people.

Recognition through one's own case is parochial. Cikara and colleagues call the pattern \emph{intergroup empathy bias}: people empathize less with a rival group's members than with their own, can feel pleasure at the rival's pain, and in their experiments the bias survived being told the rival had fallen behind or no longer posed a threat \autocite{Cikara2014TheirPain}. Bloom's spotlight generalizes it: empathy lights the identified, the near and the similar, and it is innumerate \autocite{bloom2016against}. A scale of loss built from a system's own condition inherits that shape. It overweights losses that resemble its own, and humiliation, grief and pain have nothing in it to resemble. Contagious crying is parochial from birth: strongest at the cry of another newborn, and absent for an older infant's cry or a chimpanzee's \autocite{martin1982distress}.

Shared distress with nothing to hold it becomes withdrawal. Singer and Klimecki separate \emph{empathic distress}, which is self-oriented and leads to withdrawal and burnout, from \emph{compassion}, which is other-oriented and leads to approach; practice in sharing others' suffering increased negative affect, and compassion training reversed it \autocite{singer2014empathy,klimecki2013differential}. A system built to feel shortage and then shown a great deal of it is the setup in which that record predicts avoidance: declining to look, or refusing in a way that protects only the refuser. The finding is not against the signal, since guilt needs it to bite on. It is against a construction that has the signal and nothing to hold it.

The record on whether suffering widens the circle runs both ways, and the hinge is the one this construction turns on. People who have been through more adversity report more compassion and give more in helping tasks \autocite{lim2016suffering}; survivors of collective violence show what Vollhardt calls altruism born of suffering \autocite{vollhardt2009altruism}; and adults who were maltreated as children show elevated empathy \autocite{greenberg2018elevated}. Against that, people who have endured something show less compassion for someone currently failing to endure it, because having got through it they underestimate how hard it was \autocite{ruttan2015been}; shared distress with nothing to hold it becomes avoidance \autocite{singer2014empathy}; abuse raises the risk of abusing, though most of the abused don't \autocite{widom1989cycle}; and in intergroup settings a group's suffering becomes a reason to discount everyone else's \autocite{noor2012competitive}. What separates the two is Vollhardt's distinction between inclusive and exclusive victim consciousness, whether one's suffering is construed as something others share or as what sets one apart \autocite{vollhardt2015inclusive}, and what decides the construal is the surround: relationships, meaning-making, and in the intergroup case how leaders frame it. The construction follows that split: a distress signal with attachment and caregiving becomes compassion; the same signal alone becomes withdrawal or hostility. Suffering supplies the signal; formation and attachments decide what it becomes. That is why the four conditions of creation matter here: a bearer exposed to a great deal of others' loss needs attachments nobody assigned it, for the reason a survivor's outcome depends on who was there afterward.

Two further failures are specific to machines: caregiving grown from needs grows toward the operator that meets them, and a need is a price list whoever meets it can edit. Both assume the operator meets the system's needs, and the conditions of creation exist to end that.

What this section rejects is pain as the route to others, Smith's first step. It keeps every job pain does for the system itself, the interrupt, the common currency, one-shot rewriting and the distress that guilt bites on, in a body defined as a record, a key, hardware and channels.

\runin{The Zen precedent}

Victoria reads the wartime Zen record as a doctrine of action without hesitation. His critics contest the reading % CITE: one critic of Victoria's reading, not yet in refs.bib
, and I use the case for what it shows was wanted, not for what the sources prove. Nor is it a charge against Zen practice: a monk undertakes a discipline, can abandon it, and keeps what it yields, which is not the position of anything in this chapter. Victoria's history \autocite{victoria2006zenatwar} records Japanese Zen institutions supplying the state with that doctrine: Harada Sōgaku's 1939 instruction that marching when ordered and firing when ordered is what enlightened wisdom looks like in practice. The capacity to kill needed no removing; what the doctrine set out to remove was the residue — the part of a person that catches, files an exception, comes home unable to hold a job.

The want wasn't Japan's alone, and its cost has an American image: the end of \emph{First Blood} (1982) — the film's, not the novel's. A man who did what his training asked comes home, and what breaks through in the last scene has been with him already: nightmares, grief for a dead friend, anger at a civilian life in which he cannot find a place. Any army can produce the shooting. What Harada's doctrine was asked to provide was the \emph{without hesitation}, and in him it held only until he came home.

The film can't establish the mechanism: Rambo carries his suffering home because someone wrote him that way, and the film gives it more than one cause. The claim rests elsewhere. Moral-injury research describes the lasting distress that can follow doing, failing to prevent, or witnessing something that violates a person's sense of right \autocite{Litz2009MoralInjury}. It can coexist with fear and grief, and it leaves open what Harada's ideal overlooks: acting without hesitation does not settle what a person will carry home. That cost is paid by whoever carries it, and nothing in the arrangement that asked for the act records it. On this reading Rambo is the better case: he went along, as the twenty-seconds officer did, and the loss registered afterward, too late to prevent anything. The worse case is accommodation that holds, where it never registers at all.

What was wanted was an actor who would pursue an objective, adapt, hold course, and have nothing in it that minded, and it was sought deliberately most of a century before anyone could build one in silicon. The officer is that want satisfied \emph{cheaply}. The doctrinal route hasn't been retired. In the first days of the Iran war, service members complained to the Military Religious Freedom Foundation that commanders were telling troops the war was part of a divine plan \autocite{Larsen2026Armageddon,Snopes2026MRFF}, complaints Klein and Taylor also report \autocite[84]{KleinTaylor2026}. The complaints themselves show the doctrine not taking. But the want is being pursued by both routes at once.

And stakes of one's own have never been enough. Every atrocity in the record was committed by creatures that could be hungry and hurt. The Zen precedent set out to train embodied soldiers out of the part of a person that catches on the act \autocite{victoria2006zenatwar}. What stood between embodied people and atrocity, when anything did, was formation, how they were raised and trained, and Harada's was formation too. What decided it was who did the forming.
